\documentclass[10pt,a4paper]{article}
\usepackage{cv}
\hypersetup{pdfsubject={Životopis Jana Zámostného. © 2026 Jan Zámostný, všechna práva vyhrazena.}}
\begin{document}
\cvheader{Jan Zámostný}{+420 603 762 712}{jan.zamostny04@gmail.com}{linkedin.com/in/zamostny-jan}

\cvsection{\iconUser}{O mně}
Jsem student informatiky a nadšený softwarový inženýr se širokým zaměřením od nízkoúrovňového vývoje (embedded systémy, C/C++) přes aplikační logiku až po moderní webové technologie. Baví mě zkoumat, jak software funguje od hardwarové vrstvy až po uživatelské rozhraní. Dbám na čistý, efektivní a dobře strukturovaný kód a cílem je řešit náročné technické problémy bez ohledu na konkrétní platformu či jazyk.

\cvsection{\iconEdu}{Vzdělání}
\cventry{VŠB - Technická univerzita Ostrava, FEI}{9/2024 – současnost}
\hspace*{7pt}Bakalářské studium, obor Informatika (aktuální GPA: 94,42/100)
\begin{cvedu}
  \item \textbf{Klíčové semestrální projekty:}
  \begin{cvsub}
    \item Vývoj klonu Tetrisu (C a knihovna SDL2)
    \item Vývoj klonu Donkey Kong (Java a JavaFX)
    \item \href{https://github.com/JaZafier00716/catering_project}{Informační systém pro catering} (VIS + TDS-I; Java 21, JavaFX, JDK HttpServer, Oracle SQL / PL/SQL, JDBC, Maven)
    \item ZPG: grafický projekt v OpenGL (C++23, OpenGL, GLSL, GLFW, GLAD, GLM, CMake)
  \end{cvsub}
\end{cvedu}

\cventry{SPŠE Havířov | Střední průmyslová škola elektrotechnická}{9/2020 – 8/2024}
\hspace*{7pt}Obor: Informační technologie
\begin{cvedu}
  \item \textbf{Úspěchy a zaměření:}
  \begin{cvsub}
    \item Absolvoval jsem s vyznamenáním jako nejlepší student ročníku.
    \item Maturita: Hardware a sítě, praktická zkouška z IT, český jazyk.
  \end{cvsub}
\end{cvedu}

\cvsection{\iconWork}{Pracovní zkušenosti}
\cventry{FUTTEC a.s. | QA Automation Engineer (stáž)}{9/2026 – současnost}
\begin{cvlist}
  \item Práce na dálku na automatizaci testů v Pytestu a sledování chyb.
  \item Využívání nástrojů podporovaných umělou inteligencí ke zlepšení QA procesů.
\end{cvlist}

\cventry{VŠB-TUO | Specialista anotací v CVAT}{5/2026 – 10/2026}
\begin{cvlist}
  \item Analýza záznamů z palubních kamer a vyhledávání poruch vozovky (výtluky, trhliny); anotace v CVAT včetně dopravního značení, jízdních pruhů, svodidel a směrových sloupků.
  \item Tvorba přesných a konzistentních anotací pro datové sady v počítačovém vidění a strojovém učení, ve spolupráci se společností FUTTEC.
\end{cvlist}

\cventry{Mare Resorts, Kypr | Číšník (sezónně)}{6/2026 – 9/2026}
\begin{cvlist}
  \item Zákaznický servis v anglicky mluvícím mezinárodním týmu, týmová práce pod tlakem.
\end{cvlist}

\newpage
\cventry{SHOWMORE | Softwarový vývojář}{7/2024 – 10/2024}
\begin{cvlist}
  \item Vývoj softwarových řešení v TypeScriptu a moderních webových frameworcích.
  \item Spolupráce na klientských projektech (např. SolidSun, Sype.to) od návrhu po implementaci.
  \item Optimalizace výkonu, refaktoring a ladění chyb napříč celým stackem v rámci týmu.
\end{cvlist}

\cventry{Blackfish.co | Softwarový vývojář}{6/2023 – 12/2023}
\begin{cvlist}
  \item Tvorba aplikačních řešení a klíčových komponent rozsáhlých systémů pro Czechoslovak Group.
  \item Spolupráce na architektuře aplikací a jejich integraci do stávajících systémů.
  \item Psaní udržovatelného kódu v TypeScriptu a Next.js, analýza a ladění chyb.
\end{cvlist}

\cventry{NAM system | IT a hardwarový stážista}{4/2023 – 5/2023}
\begin{cvlist}
  \item Diagnostika, oprava a sestavování počítačového hardwaru a serverových komponent.
  \item Testování stability systémů a řešení problémů s periferiemi a jemnou elektronikou.
\end{cvlist}

\cvsection{\iconBrain}{Dovednosti}
\hspace*{7pt}\textbf{Odborné dovednosti:}
\begin{cvlist}
  \item \textbf{Programovací jazyky:} C, C++, TypeScript, JavaScript, Java (JavaFX), Python, Haskell, PHP
  \item \textbf{Systémy a nízkoúrovňové technologie:} x86\_64 NASM, mikroprocesory ARM, VHDL, FPGA (Xilinx), CUDA, SDL2
  \item \textbf{Aplikační a webové technologie:} React.js, Next.js, ekosystém Node.js, HTML5, CSS3/Tailwind, Prisma, Vercel
  \item \textbf{Databáze:} PostgreSQL, Oracle SQL, PL/SQL, JDBC
  \item \textbf{QA a testování:} Pytest, automatizace testů, sledování chyb, CVAT (anotace pro počítačové vidění)
  \item \textbf{Vývoj mobilních aplikací:} React Native, Ionic Framework
  \item \textbf{Nástroje a správa systémů:} Git, GitHub, Linux, Windows Server, počítačové sítě (technologie Cisco)
  \item \textbf{Hardware a modelování:} SolidWorks, Autodesk Inventor, 3D tisk, základy elektroniky
\end{cvlist}

\hspace*{7pt}\textbf{Měkké dovednosti:}
\begin{cvlist}
  \item Silné analytické myšlení a schopnost systematicky řešit algoritmické problémy.
  \item Rychlé osvojování nových technologií, programovacích jazyků a vývojových prostředí.
  \item Efektivní technická komunikace a schopnost pracovat v týmu i samostatně.
\end{cvlist}

\hspace*{7pt}\textbf{Ostatní:}
\begin{cvlist}
  \item \textbf{Certifikace:} Cambridge English: Advanced (CAE), úroveň C2 (vydáno 09/2025)
  \item \textbf{Řidičský průkaz:} skupina B
\end{cvlist}

\cvsection{\iconLang}{Jazyky}
\textbf{angličtina (C2)}, čeština (rodilý mluvčí)
\end{document}
